\chapter{Gold Standard}
\label{chap:donnees-gold}

\section{Introduction}\label{sec:c4g-intro}

Pour évaluer un système qui extrait automatiquement des informations cliniques, il faut pouvoir comparer ce qu'il produit à ce que contiennent réellement les dossiers. Avec les modèles génératifs multimodaux, cette comparaison est indispensable : leurs réponses sont souvent bien formées et médicalement vraisemblables, alors même qu'une partie de leur contenu ne se trouve pas dans le document analysé. Dans le cadre KG-EDS~\cite{kgeds}, faute de corpus annoté, ni la fiabilité des extractions ni les hallucinations ne pouvaient être mesurées précisément. Le Gold Standard présenté ici a été construit pour lever cette limite.

Sa construction a été traitée comme une étape méthodologique à part entière. Les dossiers ont été annotés à la main selon le TRACE-Sepsis Annotation Guideline présenté au chapitre~\ref{chap:extraction}, sans l'aide d'aucun modèle, puis revus par le Pr Nicholas Heming, ce qui assure à la fois l'indépendance et la pertinence clinique de la référence.

Le chapitre décrit d'abord le sous-corpus de référence et le protocole d'annotation (section~\ref{sec:c3-3-7}), puis la construction du Gold Standard (section~\ref{sec:c3-3-8}), sa validation médicale (section~\ref{sec:c3-3-9}) et sa caractérisation (section~\ref{sec:c3-3-10}). Il définit ensuite la manière dont une extraction est comparée à la référence : son usage dans l'évaluation (section~\ref{sec:c3-3-11}), les modes d'appariement (section~\ref{sec:c4-matching}) et les métriques (section~\ref{sec:c4-metriques}). La section~\ref{sec:c3-3-12} discute les limites de la démarche.

% ============================================================
\section{Sous-corpus de référence et protocole d’annotation}\label{sec:c3-3-7}

La figure~\ref{fig:c3-3-1} résume le chemin qui mène du corpus clinique au Gold Standard final. Les étapes sont détaillées dans cette section et les deux suivantes.

\begin{figure}[H]
\centering
\includegraphics[width=\textwidth]{images/fig_gold.png}
\caption{Processus de construction du Gold Standard}\label{fig:c3-3-1}
\end{figure}

\subsection{Sélection et justification du sous-corpus}\label{sec:c3-3-7-1}

L’annotation a porté sur 49 dossiers, soit environ 20\,\% du corpus et 291 pages. Le corpus ayant été réuni progressivement (section~\ref{sec:c3-3-2-3}), les dossiers ont été retenus dans l’ordre où l’AP-HP les a mis à disposition. Ils proviennent de deux lots successifs :

\begin{itemize}
  \item \textbf{un premier lot}, disponible dès le début du projet ;
  \item \textbf{un second lot}, mis à disposition par l’AP-HP fin juillet 2026.
\end{itemize}

Une partie de chaque lot a été annotée. Commencer par le premier lot a permis de lancer l’annotation et de stabiliser le guideline sans attendre la livraison suivante ; des dossiers du second lot ont été ajoutés à mesure qu’ils arrivaient. Les dossiers n’ont donc pas été choisis pour leur facilité ou leur difficulté. Rien ne garantit en revanche que le sous-corpus soit représentatif des 246 dossiers, puisque la composition de chaque lot dépend de facteurs administratifs comme la période de prise en charge ou l’ordre d’archivage.

Annoter un sous-corpus plutôt que l’ensemble répond au coût de l’annotation experte. Chaque dossier doit être annoté selon le schéma, ses relations vérifiées et ses annotations consolidées avant la revue clinique. Mieux vaut une référence soignée sur une partie du corpus qu’une annotation superficielle de tous les dossiers. La représentativité d’un corpus annoté dépend d’ailleurs davantage de la diversité des données et de la manière de les sélectionner que de leur nombre~\cite{biber1993}, et plusieurs ressources cliniques annotées ont été construites de la même façon, à partir de sous-ensembles~\cite{roberts2007clef,patel2018annotation}.

\subsection{Annotateurs et environnement d’annotation}\label{sec:c3-3-7-2}

Les 49 dossiers ont été annotés par l’auteure de ce mémoire, en suivant le TRACE-Sepsis Annotation Guideline. La campagne a duré près de 25 jours. Trois raisons l’expliquent : le volume (291 pages), la densité des comptes rendus de réanimation, et le fait que l’annotatrice n’est pas médecin. Chaque décision devait être confrontée au guideline, et les cas douteux étaient mis de côté pour la validation médicale.

L’annotation a été faite dans Label Studio~\cite{labelstudio}, qui permet de délimiter les entités dans le document et de tracer les relations entre elles. Un outil dédié garantit une structure homogène des annotations et en facilite l’export.

L’annotation initiale a été faite entièrement à la main, sans pré-annotation et sans recours à un modèle de langage ou à un modèle multimodal. La référence est ainsi indépendante des systèmes évalués (section~\ref{sec:c3-3-10-1}). En revanche, une annotation produite par une personne qui n’est pas médecin devait être validée cliniquement : c’est l’objet de la section~\ref{sec:c3-3-9}.

\subsection{Organisation des tâches}\label{sec:c3-3-7-3}

Les dossiers étant de longueur variable, chaque page a constitué une tâche d’annotation. Ce découpage rend les documents plus lisibles, facilite la localisation des preuves et rend l’annotation plus maniable. Il ne fait pas de chaque page une référence indépendante : \textbf{la page est l’unité pratique d’annotation, le dossier l’unité finale du Gold Standard}. On conserve ainsi la précision locale des annotations sans perdre la cohérence clinique du dossier.

\subsection{Annotation des entités et des relations}\label{sec:c3-3-7-4}

Sur chaque page, les mentions relevant des catégories du schéma ont été repérées selon le guideline, en tenant compte du type de concept et du contexte. Une information insuffisamment étayée était écartée ou signalée pour vérification. Les relations ont ensuite été tracées entre les entités repérées : chacune devait respecter les signatures du guideline et s’appuyer sur un lien écrit ou documenté, jamais sur la seule proximité des entités.

\subsection{Contrôle des annotations}\label{sec:c3-3-7-5}

Avant la validation médicale, un contrôle a permis de repérer les anomalies qui auraient dégradé la référence :

\begin{itemize}
  \item frontières incorrectes et mauvais types ;
  \item doublons et chevauchements injustifiés ;
  \item relations dont une extrémité manque ;
  \item signatures de relation invalides ;
  \item annotations créées malgré une négation ;
  \item informations reposant sur une inférence non justifiée.
\end{itemize}

\section{Construction du Gold Standard}\label{sec:c3-3-8}

\subsection{Consolidation à l’échelle du dossier}\label{sec:c3-3-8-1}

Les annotations faites page par page ont ensuite été réunies pour former une représentation de chaque dossier. Les concaténer n’aurait pas suffi, car plusieurs pages d’un même dossier peuvent mentionner le même concept, de façon répétée ou complémentaire. La consolidation a consisté à regrouper les annotations d’un même patient, à harmoniser les mentions équivalentes, à retirer les répétitions inutiles et à garder la provenance de chaque information.

Toutes les répétitions ne disparaissent pas. Une information qui a une portée temporelle, clinique ou contextuelle différente reste représentée séparément : deux dosages d’un même biomarqueur à des moments différents ne sont pas fusionnés si leur évolution a un intérêt clinique. La consolidation garde ainsi la richesse du dossier sans multiplier artificiellement des informations identiques.

\subsection{Contrôle structurel}\label{sec:c3-3-8-2}

Un contrôle structurel a ensuite vérifié que les annotations respectaient le schéma TRACE-Sepsis : types d’entités valides, relations appartenant au vocabulaire contrôlé, cohérence entre sujet, relation et objet, présence réelle des entités utilisées dans les relations, absence d’éléments techniques ou de catégories non résolues. Une relation n’a été gardée que si son sujet et son objet relevaient du domaine et de la portée prévus. Les incohérences formelles ont ainsi été éliminées avant la revue clinique.

\subsection{Mise en conformité avec la version 2.0}\label{sec:c3-3-8-3}

L’annotation ayant commencé avec la première version du guideline, elle a été mise en conformité avec la version 2.0 par des programmes qui appliquent uniquement les règles écrites du guideline. Ils ont rattaché chaque foyer infectieux au diagnostic dont il est la localisation, retiré les foyers qui localisaient un antécédent, retypé en \entite{LABEL\_NOSOLOGIQUE} les diagnostics de la liste annotés comme foyers, et supprimé les relations devenues hors signature. Ces programmes ne lisent aucune sortie de modèle. Les cas qu’ils ne pouvaient pas trancher seuls, par exemple un foyer pouvant se rattacher à plusieurs diagnostics, ont été listés puis examinés un par un en appliquant le guideline.

\subsection{Format de représentation et traçabilité}\label{sec:c3-3-8-4}

L’export de Label Studio conserve les annotations page par page : entités, types, positions, et relations définies par les identifiants des entités. Un post-traitement regroupe ensuite les annotations d’un même dossier, résout les références entre entités et relations, harmonise les mentions équivalentes et produit une représentation par dossier patient.

Le Gold Standard consolidé comporte deux parties. Les entités cliniques sont accompagnées de ce qu’il faut pour les identifier et retrouver leur origine ; les relations sont des liens orientés entre elles. Chaque information garde, quand c’est utile, sa page d’origine et une preuve. Cette preuve est volontairement courte, limitée au fragment qui justifie l’annotation, pour être vérifiée rapidement sans recopier de longs passages du dossier. Cette traçabilité sert pendant la validation médicale, puis lors de l’analyse des écarts entre le Gold Standard et les extractions des modèles.

On obtient à ce stade une première version structurée et contrôlée de la référence : le Gold Standard initial.

\section{Validation médicale du Gold Standard}\label{sec:c3-3-9}

Une annotation peut être structurellement correcte et cliniquement fausse : une relation peut respecter toutes les contraintes du schéma sans correspondre à ce que décrit le dossier. Après la validation du guideline (section~\ref{sec:c3-3-3-4}), une seconde validation médicale a donc porté sur les annotations elles-mêmes. Le Pr Nicholas Heming l’a menée sur le Gold Standard initial, après l’annotation et la consolidation. Elle a demandé en moyenne environ 2,5 heures de travail par dossier avec le médecin.

Le tableau~\ref{tab:c3-3-val} résume les deux niveaux de validation médicale.

\begin{table}[H]
\centering
\caption{Les deux niveaux de validation médicale}\label{tab:c3-3-val}
\small
\begin{tabular}{>{\raggedright\arraybackslash}p{0.17\linewidth}>{\raggedright\arraybackslash}p{0.30\linewidth}>{\raggedright\arraybackslash}p{0.17\linewidth}>{\raggedright\arraybackslash}p{0.22\linewidth}}
\toprule
\textbf{Objet} & \textbf{Contenu examiné} & \textbf{Durée} & \textbf{Moment} \\
\midrule
Guideline & Catégories d’entités, vocabulaires contrôlés, relations autorisées, règles d’annotation & Environ 5 heures & Avant la campagne d’annotation \\
Annotations & Entités et relations du Gold Standard initial & Environ 2,5 heures par dossier en moyenne & Après annotation et consolidation \\
\bottomrule
\end{tabular}
\end{table}

\subsection{Objectif et protocole}\label{sec:c3-3-9-1}

Il s’agit de vérifier que la représentation structurée reste fidèle à ce que le dossier documente. Le médecin a confronté les annotations à leur source, présentées avec des preuves aussi courtes et ciblées que possible. Le but n’était ni de reconstituer le dossier ni de produire une seconde annotation indépendante, mais de faire une revue experte du Gold Standard initial pour savoir quelles annotations confirmer et lesquelles corriger.

\subsection{Validation des entités}\label{sec:c3-3-9-2}

Pour chaque entité examinée, la revue a porté sur sa présence réelle dans le dossier, l’adéquation entre la mention et le type attribué, la délimitation de la mention, son interprétation dans le contexte clinique, et la prise en compte d’une négation ou d’une incertitude éventuelles. Les catégories voisines demandaient une attention particulière : une manifestation clinique doit être distinguée d’un diagnostic ou d’une défaillance d’organe, et une pathologie préexistante d’un diagnostic lié à l’épisode en cours.

\subsection{Validation des relations}\label{sec:c3-3-9-3}

Une relation doit être structurellement valide, mais aussi correspondre au lien clinique réellement décrit. La revue a porté sur l’entité source, l’entité cible, le type de relation, son sens, et sur ce qui, dans le dossier, justifie l’association. Un triplet peut en effet respecter la signature \entite{TRAITEMENT} $\rightarrow$ relation $\rightarrow$ \entite{MICRO\_ORGANISME} sans que le dossier permette de relier ce traitement à ce germe. La revue médicale fait apparaître ces relations formellement correctes mais mal étayées.

\subsection{Intégration des corrections}\label{sec:c3-3-9-4}

À l’issue de la revue, chaque annotation a été confirmée, modifiée (information présente mais mal représentée), supprimée (justification insuffisante) ou ajoutée (information pertinente oubliée lors de l’annotation initiale), en restant cohérent avec le guideline. Lorsqu’une correction du médecin traduisait une règle générale, cette règle a été appliquée à tous les dossiers concernés, et pas seulement au dossier corrigé. C’est le cas, par exemple, de l’exclusion des noms d’établissements des \entite{SERVICE\_MEDICAL}, ou de la non-annotation des valeurs témoins du laboratoire. Lorsqu’une correction contredisait une règle écrite du guideline, c’est la règle du guideline qui a été appliquée, afin que le référentiel reste cohérent d’un dossier à l’autre.

\section{Gold Standard final et caractérisation}\label{sec:c3-3-10}

\subsection{Version finale et indépendance de la référence}\label{sec:c3-3-10-1}

Une fois les corrections intégrées, une dernière consolidation a vérifié qu’elles n’introduisaient pas de nouvelles incohérences et que les entités et les relations respectaient toujours le schéma TRACE-Sepsis v2.0. Le résultat est le Gold Standard final, utilisé comme référence dans les expériences.

Cette référence a été construite sans les modèles multimodaux évalués : ni les prédictions de Qwen2.5-VL, ni celles de Mistral Small 4, ni celles d’un autre modèle n’en déterminent le contenu. On évite ainsi un raisonnement circulaire où un modèle contribuerait à définir sa propre vérité de référence. L’indépendance doit valoir dans les deux sens : aucune sortie évaluée ne doit avoir été construite ou corrigée à l’aide du Gold Standard. Dans la chaîne de traitement, seule l’étape d’évaluation a accès à la référence (chapitre~\ref{chap:architecture}).

\subsection{Statistiques descriptives}\label{sec:c3-3-10-2}

Le tableau~\ref{tab:c3-3-4} donne les principales statistiques du Gold Standard final, compté au niveau des annotations de pages. Les tableaux~\ref{tab:c3-3-5} et~\ref{tab:c3-3-6} détaillent la répartition par type d’entité et par type de relation.

\begin{table}[H]
\centering
\caption{Statistiques globales du Gold Standard final (49 dossiers, 291 pages)}\label{tab:c3-3-4}
\small
\begin{tabular}{>{\raggedright\arraybackslash}p{0.346\linewidth}>{\raggedright\arraybackslash}p{0.260\linewidth}>{\raggedright\arraybackslash}p{0.260\linewidth}}
\toprule
\textbf{Indicateur} & \textbf{Entités} & \textbf{Relations} \\
\midrule
Total & 12\,174 & 5\,543 \\
Moyenne par dossier (écart-type) & 248,4 (82,2) & 113,1 (46,6) \\
Minimum -- maximum par dossier & 90 -- 474 & 25 -- 242 \\
\bottomrule
\end{tabular}
\end{table}

\begin{table}[H]
\centering
\caption{Répartition des entités par type}\label{tab:c3-3-5}
\small
\begin{tabular}{>{\raggedright\arraybackslash}p{0.432\linewidth}>{\raggedright\arraybackslash}p{0.216\linewidth}>{\raggedright\arraybackslash}p{0.216\linewidth}}
\toprule
\textbf{Type d’entité} & \textbf{Effectif} & \textbf{Proportion (\%)} \\
\midrule
\entite{BIOMARQUEUR} & 2\,989 & 24,6 \\
\entite{TRAITEMENT} & 2\,465 & 20,2 \\
\entite{SYMPTOME} & 1\,177 & 9,7 \\
\entite{POSOLOGIE} & 1\,056 & 8,7 \\
\entite{COMORBIDITE\_\allowbreak{}ANTECEDENT} & 585 & 4,8 \\
\entite{SIGNE\_\allowbreak{}VITAL} & 565 & 4,6 \\
\entite{LABEL\_\allowbreak{}NOSOLOGIQUE} & 500 & 4,1 \\
\entite{DEFAILLANCE\_\allowbreak{}ORGANE} & 485 & 4,0 \\
\entite{IMAGERIE\_\allowbreak{}PROCEDURE} & 470 & 3,9 \\
\entite{DONNEE\_\allowbreak{}PATIENT} & 438 & 3,6 \\
\entite{EVENEMENT\_\allowbreak{}TEMPOREL} & 403 & 3,3 \\
\entite{EVOLUTION\_\allowbreak{}PRONOSTIC} & 240 & 2,0 \\
\entite{MICRO\_\allowbreak{}ORGANISME} & 235 & 1,9 \\
\entite{FOYER\_\allowbreak{}INFECTIEUX} & 187 & 1,5 \\
\entite{SERVICE\_\allowbreak{}MEDICAL} & 184 & 1,5 \\
\entite{CONTEXTE\_\allowbreak{}ACQUISITION} & 125 & 1,0 \\
\entite{SCORE\_\allowbreak{}NEUROLOGIQUE} & 67 & 0,6 \\
\entite{SCORE\_\allowbreak{}SOFA} & 3 & 0,0 \\
\midrule
\textbf{Total} & 12\,174 & 100,0 \\
\bottomrule
\end{tabular}
\end{table}

\begingroup\small
\begin{longtable}{>{\raggedright\arraybackslash}p{0.200\linewidth}>{\raggedright\arraybackslash}p{0.380\linewidth}>{\raggedright\arraybackslash}p{0.120\linewidth}>{\raggedright\arraybackslash}p{0.140\linewidth}}
\caption{Répartition des relations par famille et par type}\label{tab:c3-3-6}\\
\toprule
\textbf{Famille} & \textbf{Type de relation} & \textbf{Effectif} & \textbf{Proportion (\%)} \\
\midrule
\endfirsthead
\multicolumn{4}{l}{\emph{Tableau~\ref{tab:c3-3-6} (suite)}}\\
\toprule
\textbf{Famille} & \textbf{Type de relation} & \textbf{Effectif} & \textbf{Proportion (\%)} \\
\midrule
\endhead
\midrule
\multicolumn{4}{r}{\emph{Suite page suivante}}\\
\endfoot
\bottomrule
\endlastfoot
Thérapeutiques & traitement\_\allowbreak{}administre\_\allowbreak{}a & 1\,861 & 33,6 \\
 & traitement\_\allowbreak{}a\_\allowbreak{}pour\_\allowbreak{}posologie & 1\,012 & 18,3 \\
 & traitement\_\allowbreak{}cible\_\allowbreak{}defaillance & 65 & 1,2 \\
 & traitement\_\allowbreak{}indique\_\allowbreak{}par\_\allowbreak{}label\_\allowbreak{}nosologique & 58 & 1,0 \\
 & antibiotique\_\allowbreak{}adequat\_\allowbreak{}pour & 37 & 0,7 \\
 & antibiotique\_\allowbreak{}inadequat\_\allowbreak{}pour & 8 & 0,1 \\
\addlinespace
Symptômes & presente\_\allowbreak{}symptome & 936 & 16,9 \\
\addlinespace
Diagnostiques & presente\_\allowbreak{}dysfonction\_\allowbreak{}organe & 410 & 7,4 \\
 & a\_\allowbreak{}pour\_\allowbreak{}label\_\allowbreak{}nosologique & 326 & 5,9 \\
 & label\_\allowbreak{}nosologique\_\allowbreak{}a\_\allowbreak{}pour\_\allowbreak{}foyer & 80 & 1,4 \\
 & score\_\allowbreak{}sofa\_\allowbreak{}calcule\_\allowbreak{}pour & 3 & 0,1 \\
\addlinespace
Temporelles et pronostiques & patient\_\allowbreak{}a\_\allowbreak{}pour\_\allowbreak{}evolution & 163 & 2,9 \\
 & service\_\allowbreak{}medical\_\allowbreak{}accueille & 145 & 2,6 \\
 & evenement\_\allowbreak{}associe\_\allowbreak{}evolution & 20 & 0,4 \\
 & precede & 2 & 0,0 \\
\addlinespace
Biomarqueurs et microbiologie & biomarqueur\_\allowbreak{}supporte\_\allowbreak{}defaillance\_\allowbreak{}organe & 114 & 2,1 \\
 & micro\_\allowbreak{}organisme\_\allowbreak{}isole\_\allowbreak{}dans & 98 & 1,8 \\
 & biomarqueur\_\allowbreak{}est\_\allowbreak{}critere\_\allowbreak{}de & 57 & 1,0 \\
 & biomarqueur\_\allowbreak{}obtenu\_\allowbreak{}a & 22 & 0,4 \\
\addlinespace
Imagerie & imagerie\_\allowbreak{}objective\_\allowbreak{}foyer & 69 & 1,2 \\
\addlinespace
Signes vitaux & signe\_\allowbreak{}vital\_\allowbreak{}supporte\_\allowbreak{}defaillance\_\allowbreak{}organe & 39 & 0,7 \\
\addlinespace
Terrain et risque & comorbidite\_\allowbreak{}est\_\allowbreak{}facteur\_\allowbreak{}risque\_\allowbreak{}de & 15 & 0,3 \\
 & contexte\_\allowbreak{}acquisition\_\allowbreak{}influence\_\allowbreak{}risque\_\allowbreak{}resistance & 3 & 0,1 \\
\midrule
\textbf{Total} & & 5\,543 & 100,0 \\
\end{longtable}
\endgroup

Ces répartitions sont très déséquilibrées. Les biomarqueurs et les traitements représentent à eux seuls près de la moitié des entités, et deux relations, l’administration d’un traitement et sa posologie, plus de la moitié des relations. À l’inverse, deux types d’entités (\entite{SCORE\_qSOFA} et \entite{STADE\_IRA}) et onze types de relations n’apparaissent jamais, notamment les relations liées aux scores SOFA et qSOFA, au stade d’IRA et certaines relations portant sur les symptômes et le site infecté. Ces écarts doivent être gardés en tête pour interpréter les résultats : un bon score global peut masquer des performances faibles sur des catégories rares, et aucune performance ne peut être mesurée sur une catégorie absente de la référence. Les écarts entre dossiers, de 90 à 474 entités, permettent par ailleurs de situer les performances selon la complexité des documents.

\subsection{Effet de la validation médicale}\label{sec:c3-3-10-3}

Le tableau~\ref{tab:c3-3-7} compare la référence avant et après la revue médicale. Il ne s’agit pas d’une mesure d’accord inter-annotateurs, mais d’une façon de chiffrer ce que la validation a apporté à la consolidation du Gold Standard.

\begin{table}[H]
\centering
\caption{Gold Standard avant et après validation médicale}\label{tab:c3-3-7}
\small
\begin{tabular}{>{\raggedright\arraybackslash}p{0.404\linewidth}>{\raggedright\arraybackslash}p{0.231\linewidth}>{\raggedright\arraybackslash}p{0.231\linewidth}}
\toprule
\textbf{Indicateur} & \textbf{Avant validation} & \textbf{Après validation} \\
\midrule
Nombre total d’entités & \acompleter{à compléter} & 12\,174 \\
Nombre total de relations & \acompleter{à compléter} & 5\,543 \\
Entités ajoutées / supprimées / reclassées & --- & \acompleter{à compléter} \\
Relations ajoutées / supprimées / corrigées & --- & \acompleter{à compléter} \\
\bottomrule
\end{tabular}
\end{table}

\section{Utilisation du Gold Standard dans l’évaluation}\label{sec:c3-3-11}

Le Gold Standard final sert de référence pour mesurer les performances des systèmes d’extraction. Pour chaque dossier, la sortie structurée d’un modèle est comparée à la référence, ce qui permet de compter les prédictions correctes (vrais positifs, TP), les informations en trop (faux positifs, FP) et les informations manquantes (faux négatifs, FN). La précision, le rappel et le F1-score s’en déduisent~\cite{vanrijsbergen1979} :

\begin{equation}
\text{Précision} = \frac{TP}{TP + FP}
\end{equation}
\begin{equation}
\text{Rappel} = \frac{TP}{TP + FN}
\end{equation}
\begin{equation}
F_1 = \frac{2 \times \text{Précision} \times \text{Rappel}}{\text{Précision} + \text{Rappel}}
\end{equation}

Ces mesures sont calculées séparément pour les entités et pour les relations, selon les modes d’appariement (\emph{matching}) décrits à la section~\ref{sec:c4-matching}.

Au-delà des scores globaux, la comparaison détaillée avec la référence permet de qualifier les écarts : information de la référence non retrouvée, entité en trop, mauvais type, valeur ou unité erronée, relation manquante, relation que le dossier ne justifie pas, sens inversé, incohérence structurelle, temporelle ou contextuelle. Ces écarts sont la base de l’analyse des erreurs et des hallucinations présentée dans la suite.

Le Gold Standard ne change pas pendant l’évaluation de la chaîne TRACE : les sorties de chaque étape sont comparées à la même référence. On peut ainsi savoir si une étape corrige réellement une extraction ou si elle se contente d’en changer la forme.

\section{Appariement entre Gold Standard et extraction}\label{sec:c4-matching}

Pour comparer une extraction au Gold Standard, il faut décider quand un élément prédit correspond à un élément de référence. Une correspondance exacte est trop sévère pour un modèle génératif, qui reformule souvent légèrement une mention correcte~\cite{wadhwa2023revisiting} ; une correspondance trop souple compterait comme justes des éléments erronés. Nous utilisons donc trois modes d'appariement de sévérité croissante, appliqués de la même façon à toutes les sorties évaluées dans ce mémoire, avant TRACE comme après chaque étage.

\subsection{Normalisation des textes}

Avant toute comparaison, les mentions sont normalisées : passage en minuscules, suppression des accents, unification des apostrophes, remplacement des tirets bas par des espaces, remplacement de la virgule décimale par un point (\og 37,2 \fg{} devient \og 37.2 \fg{}) et réduction des espaces multiples. Pour les entités, les apostrophes, la ponctuation interne et les tirets sont en outre remplacés par des espaces. Deux mentions qui ne diffèrent que par ces détails de forme sont ainsi considérées comme identiques.

\subsection{Mesures de proximité entre deux mentions}

Pour les entités, la similarité de deux mentions est le ratio de Ratcliff et Obershelp~\cite{ratcliff1988}, qui rapporte le nombre de caractères appariés à la longueur totale des deux chaînes :
\begin{equation}
\mathrm{sim}_E(a, b) = \frac{2M}{|a| + |b|},
\end{equation}
où $M$ est le nombre de caractères communs trouvés en recherchant récursivement les plus longues sous-chaînes communes, et $|a|$ la longueur de $a$.

Pour les relations, les extrémités sont comparées par la similarité de Levenshtein~\cite{levenshtein1966} :
\begin{equation}
\mathrm{sim}_R(a, b) = 1 - \frac{d_{\mathrm{Lev}}(a, b)}{\max(|a|, |b|)}
\end{equation}
où $d_{\mathrm{Lev}}(a,b)$ est le nombre minimal d'insertions, de suppressions et de substitutions de caractères nécessaires pour transformer $a$ en $b$. La division par la longueur de la plus longue chaîne ramène la distance entre 0 et 1, ce qui est la normalisation habituelle de la distance d'édition~\cite{navarro2001}.

Le score d'inclusion vérifie si la mention la plus courte apparaît, comme suite de mots complets, dans la plus longue. Il traduit une idée classique de l'évaluation de la reconnaissance d'entités biomédicales : une mention dont seules les bornes diffèrent de la référence peut être comptée comme correcte selon un critère moins strict que la correspondance exacte~\cite{tsai2006criteria}. Nous le définissons ainsi :
\begin{equation}
\mathrm{inc}(a, b) =
\begin{cases}
\dfrac{|\text{courte}|}{|\text{longue}|} & \text{si la courte est incluse dans la longue et } |\text{courte}| \geq 4,\\[2mm]
0 & \text{sinon.}
\end{cases}
\end{equation}
La longueur minimale de quatre caractères évite qu'une abréviation très courte soit retrouvée par hasard dans une mention plus longue. Pour les entités, une inclusion qui couvre moins de 35~\% de la mention la plus longue n'est pas retenue. Ces deux seuils, comme ceux des correspondances approximatives, ont été fixés empiriquement pour ce travail.

\subsection{Appariement des entités}

Une entité prédite ne peut être appariée qu'à une entité de référence située sur la même page et de même type. Les appariements sont recherchés en trois passes successives, chaque élément ne pouvant être apparié qu'une seule fois :
\begin{enumerate}
  \item \textbf{correspondance exacte} : les deux mentions normalisées sont identiques ;
  \item \textbf{correspondance par inclusion} : $\mathrm{inc}(a,b) > 0$, les meilleurs scores étant retenus en premier (par exemple \og noradrénaline \fg{} et \og noradrénaline IVSE \fg{}) ;
  \item \textbf{correspondance approximative} : $\mathrm{sim}_E(a,b) \geq 0{,}85$, seuil fixé empiriquement, les meilleurs scores étant retenus en premier.
\end{enumerate}

\subsection{Appariement des relations}

Les relations sont comparées à l'échelle du dossier. Le Gold Standard étant consolidé par dossier, un même triplet prédit sur plusieurs pages (par exemple un traitement relié au patient sur chaque page) n'est compté qu'une fois. Une relation prédite ne peut être appariée qu'à une relation de référence de même type. Ses deux extrémités sont comparées par leur texte, à l'aide du score
\begin{equation}
e(a, b) = \max\big(\mathrm{sim}_R(a,b),\ \mathrm{inc}(a,b)\big),
\end{equation}
et la proximité de deux relations est la moyenne des scores de leur sujet et de leur objet. Cette combinaison est propre à notre évaluation ; elle applique aux relations, extrémité par extrémité, la tolérance aux variations de surface recommandée pour évaluer l'extraction de relations par un modèle génératif~\cite{wadhwa2023revisiting} :
\begin{equation}
s(r_g, r_p) = \frac{e(\mathrm{sujet}_g, \mathrm{sujet}_p) + e(\mathrm{objet}_g, \mathrm{objet}_p)}{2}.
\end{equation}
Les trois passes sont les mêmes que pour les entités : correspondance exacte du sujet et de l'objet normalisés ; correspondance par inclusion, lorsque chaque extrémité est identique ou incluse dans celle de la référence ; correspondance approximative, lorsque $s(r_g, r_p) \geq 0{,}80$.

\subsection{Trois modes d'appariement}

Les trois passes définissent trois modes d'évaluation emboîtés, sur le modèle des critères exact, de bornes et approché distingués pour la reconnaissance d'entités biomédicales~\cite{tsai2006criteria} :
\begin{itemize}
  \item \textbf{Strict} : seules les correspondances exactes sont comptées comme correctes ;
  \item \textbf{Boundary-Relaxed} : les correspondances exactes et par inclusion sont comptées comme correctes, ce qui tolère une différence de délimitation de la mention ;
  \item \textbf{Relaxed} : les trois types de correspondance sont comptés comme corrects, ce qui tolère en plus de légères variations d'écriture.
\end{itemize}
Dans tous les modes, un élément retrouvé avec le bon texte mais un type différent (une entité de mauvais type, ou une relation de mauvais type entre les mêmes extrémités) est compté comme une erreur, c'est-à-dire comme un faux positif et un faux négatif. L'indicateur \emph{nie} n'intervient pas dans l'appariement : la négation est portée par la mention elle-même (\og pas de toux \fg{}). Sauf mention contraire, les résultats de ce mémoire sont donnés en mode Relaxed.

% ------------------------------------------------------------
\section{Métriques d'évaluation}\label{sec:c4-metriques}

À partir des appariements, on compte pour chaque dossier les vrais positifs, les faux positifs et les faux négatifs, puis on calcule la précision, le rappel et le F1-score définis à la section~\ref{sec:c3-3-11}~\cite{vanrijsbergen1979}. Ces mesures sont agrégées de deux façons~\cite{sokolova2009} :
\begin{itemize}
  \item la \textbf{micro-moyenne} additionne les vrais positifs, faux positifs et faux négatifs de tous les dossiers avant de calculer les mesures ; chaque élément pèse le même poids, si bien que les dossiers les plus riches pèsent davantage ;
  \item la \textbf{macro-moyenne} calcule le F1-score de chaque dossier, puis en fait la moyenne ; chaque dossier pèse alors le même poids, quelle que soit sa taille.
\end{itemize}
Les deux agrégations sont complémentaires : un écart important entre elles signale que les performances varient fortement d'un dossier à l'autre.

Nous rapportons enfin le taux d'hallucination, défini comme la part des éléments extraits qui ne correspondent à aucun élément de la référence. Cette définition reprend celle de la métrique CHAIR, qui mesure en description d'images la proportion d'objets mentionnés par le modèle et absents de la vérité terrain~\cite{rohrbach2018chair} :
\begin{equation}
H = 1 - \text{Précision} = \frac{FP}{TP + FP}.
\end{equation}
Ce taux mesure la proportion d'informations produites par le modèle que le Gold Standard ne justifie pas. Il ne distingue pas, à lui seul, une information inventée d'une information correcte mais absente de la référence ; cette distinction relève de l'analyse des erreurs.

% ------------------------------------------------------------
\section{Limites}\label{sec:c3-3-12}

Plusieurs limites sont à garder en tête pour interpréter les résultats obtenus avec cette référence.

\begin{itemize}
  \item \textbf{Annotation.} Une seule personne a annoté le corpus et aucun accord inter-annotateurs n’a été mesuré. Une double annotation d’un échantillon de dossiers, avec un coefficient kappa~\cite{cohen1960} ou un F1 inter-annotateurs, permettrait d’estimer la reproductibilité du guideline.
  \item \textbf{Validation médicale.} Un seul médecin, le Pr Heming, a validé le guideline et les annotations. Il s’agit d’une revue experte destinée à fiabiliser la référence, et non d’une étude de concordance entre plusieurs évaluateurs.
  \item \textbf{Taille et représentativité.} Les 49 dossiers annotés (291 pages) représentent environ 20\,\% des 246 dossiers. Issus de deux lots successifs et choisis selon leur ordre de mise à disposition, et non par tirage aléatoire ou stratifié, ils ne forment pas forcément un échantillon représentatif. Certaines catégories y sont rares ou absentes (section~\ref{sec:c3-3-10-2}), ce qui limite la portée des résultats par catégorie.
  \item \textbf{Pseudonymisation.} Le masquage au stylo fait avant la numérisation ne couvre qu’une partie des noms, et des informations indirectement identifiantes restent visibles dans les images. La pseudonymisation automatique appliquée ensuite porte sur les textes transcrits et les annotations, pas sur les images des pages.
  \item \textbf{Spécificité du domaine.} Le guideline a été conçu pour le sepsis ; l’appliquer à d’autres contextes cliniques demanderait d’adapter le schéma et les règles.
\end{itemize}

\section{Conclusion}\label{sec:c3-3-13}

Ce chapitre a décrit la construction de la référence clinique qui sert à évaluer l'extraction et TRACE. Un sous-corpus de 49 dossiers (291 pages) a été annoté à la main dans Label Studio selon le TRACE-Sepsis Annotation Guideline, avec la page comme unité de travail et le dossier comme unité de référence. Les annotations ont été contrôlées, consolidées, mises en conformité avec la version~2.0 du guideline, puis revues par le Pr Nicholas Heming ; ses corrections ont été intégrées avant une dernière consolidation, qui a donné le Gold Standard final.

Le chapitre a aussi défini la manière dont une extraction est comparée à cette référence, avec trois modes d'appariement et des métriques en micro- et macro-moyenne, qui servent à l'évaluation de tout le mémoire. Indépendante des modèles évalués et traçable jusqu'aux documents, cette référence permet de mesurer les performances des systèmes d'extraction et d'analyser les écarts entre information générée et information de référence. Le chapitre suivant présente TRACE, la chaîne de vérification construite pour détecter et réduire les hallucinations.
